%% kapa: exam=1 round=19 year=2008 period=2 form=A subject=english subject_ko=영어 seq=30 number=70
%% source: Q-Net 감정평가사 기출문제 (cst00309) artlSeq=1001962 file=extracted/1008/1001962/2008년 제19회 감정평가사 1차시험 기출문제/2교시.hwp
%% conversion: kordoc 4.24.51 -> tools/gen1.py (auto); review=proofread-sonnet; flags=none
\begin{question}{70}
다음 글의 제목으로 가장 알맞은 것을 고르시오.
\begin{framedtext}
Perhaps the greatest influence that Protestantism has had on American life comes from its philosophy regarding a person's relationship to work. This philosophy――commonly called the Protestant work ethic――stresses the moral value of work, self-discipline, and personal responsibility. According to this ethic, people prove their worth to themselves and to God by working hard, being honest and thrifty, and avoiding luxury, excessive pleasure, and waste. The accumulation of wealth is not considered evil unless it leads to a life of idleness and sin. The Protestant work ethic has much in common with capitalism and with the American emphasis on financial success, practicality, efficiency, and self-sufficiency.
\end{framedtext}
\begin{choices}
  \item Protestantism and Ecology
  \item Churchgoers and Their Responsibility
  \item American Moral Crisis
  \item American Capitalism
  \item Protestant Work Ethic
\end{choices}
\end{question}
\answer{5}
